\chapter{Profissionalismo e disciplina}
\label{cap:profissional}

\proximasaulas

Saber como a JVM funciona só vale se virar decisão melhor no dia a dia:
o que prometer, quando otimizar, como manter o código e onde deixar o
framework. Este capítulo resume as bases que o professor pediu, os livros
de Robert C. Martin (\citetitle{livro:martin2011}, \citetitle{livro:martin2008}
e \citetitle{livro:martin2017}), e as liga ao resto da Parte~II e ao
Quarkus. Os resumos seguem a estrutura dos livros; leia os capítulos
indicados para os argumentos completos.

\section{\emph{The Clean Coder}: o profissional}

O \citetitle{livro:martin2011} não é um livro de código; é sobre a
conduta de quem escreve código. Os capítulos mais úteis para a turma:

\begin{description}
\item[Professionalism (cap.~1)] Seções \enquote{Taking Responsibility},
  \enquote{First, Do No Harm} e \enquote{Work Ethic}: responder pelo que
  você entrega, inclusive pelos defeitos, e cuidar da própria formação.
\item[Saying No e Saying Yes (caps.~2 e 3)] \enquote{The Cost of Saying
  Yes} e \enquote{A Language of Commitment}: um \enquote{sim} é um
  compromisso, e negociar o escopo antes custa menos do que prometer e
  não cumprir.
\item[Practicing (cap.~6)] Treinar fora do horário de entrega, como os
  katas da Parte~I.
\item[Estimation e Pressure (caps.~10 e 11)] O capítulo de estimativa
  pergunta \enquote{What Is an Estimate?} e apresenta o PERT, com três
  números em vez de um; o de pressão trata de evitá-la e de lidar com
  ela sem abandonar as disciplinas.
\end{description}

\paragraph{Ligação com a JVM.}
\enquote{Está lento, é o GC} ou \enquote{reescreve com \texttt{StringBuilder}
que fica mais rápido} são promessas sem medição. A postura profissional
que o livro defende pede que a afirmação venha com dado: um perfil do
JFR, uma listagem do \texttt{javap}, um benchmark revisado. Os
capítulos~\ref{cap:jvm-por-dentro} e~\ref{cap:lendo-bytecode} dão as
ferramentas para isso.

\section{\emph{Clean Code} e \emph{Effective Java}: o código}

O \citetitle{livro:martin2008} trata do código em si: nomes (cap.~2),
funções (cap.~3), comentários (cap.~4), tratamento de erros (cap.~7),
testes (cap.~9), classes (cap.~10), concorrência (cap.~13) e um catálogo
de \termo{code smells}{\enquote{Cheiros} de código: sinais de que algo
está mal estruturado, mesmo que funcione.} (cap.~17). O
\citetitle{livro:bloch2017} complementa com conselhos específicos de
Java, organizados em itens numerados. Vários deles ganham explicação
concreta no bytecode desta parte:

\small
\begin{longtable}{@{}p{5.6cm} p{6.6cm}@{}}
\toprule
Item do \emph{Effective Java} (3ª ed.) & Onde o bytecode mostra \\
\midrule \endhead
6. Avoid creating unnecessary objects & Autoboxing escondido,
  seção~\ref{sec:autoboxing} \\
9. Prefer try-with-resources to try-finally & \texttt{addSuppressed}
  gerado pelo \texttt{javac}, seção~\ref{sec:twr} \\
42. Prefer lambdas to anonymous classes & \texttt{invokedynamic} em vez
  de uma classe a mais no JAR, seção~\ref{sec:invokedynamic} \\
61. Prefer primitive types to boxed primitives & \texttt{==} entre
  \texttt{Integer}, seção~\ref{sec:autoboxing} \\
63. Beware the performance of string concatenation & Concatenação via
  \texttt{invokedynamic}, seção~\ref{sec:invokedynamic} \\
67. Optimize judiciously & JIT, warm-up e medição,
  capítulo~\ref{cap:jvm-por-dentro} \\
78. Synchronize access to shared mutable data & \texttt{contador++} são
  três instruções, seção~\ref{sec:modelo-memoria} \\
\bottomrule
\end{longtable}
\normalsize

\paragraph{Código limpo e desempenho não brigam (quase nunca).}
A JVM otimiza o código que roda muito, no momento em que ele roda
(capítulo~\ref{cap:jvm-por-dentro}). Métodos pequenos, que o
\emph{Clean Code} recomenda, são candidatos naturais a
\termo{inlining}{Otimização do compilador JIT (que compila o código
durante a execução; ver seção~\ref{sec:jit}) que copia o corpo do método
chamado para dentro do chamador, eliminando a chamada.}. Antes de sacrificar
legibilidade por desempenho, meça; o item 67 do \emph{Effective Java}
diz o mesmo.

\section{\emph{Clean Architecture}: onde fica o framework}

O \citetitle{livro:martin2017} organiza-se em princípios de design
(SOLID, caps.~7 a 11), princípios de componentes (caps.~12 a 14) e
arquitetura (caps.~15 a 29). Dois capítulos importam para quem usa
Quarkus todo dia:

\begin{description}
\item[The Clean Architecture (cap.~22)] Apresenta a \emph{Dependency
  Rule}: as dependências do código-fonte apontam para dentro, das
  camadas de detalhe (web, banco, framework) para as regras de negócio,
  nunca o contrário.
\item[Frameworks Are Details (cap.~32)] O framework é um detalhe; o
  livro fala de um \enquote{casamento assimétrico} (\emph{asymmetric
  marriage}) em que você se compromete com o framework e ele não se
  compromete com você.
\end{description}

\paragraph{Ponte com Quarkus.}
O Quarkus faz muito trabalho em tempo de build: o guia
\citetitle{doc:quarkus-cdi-reference} diz que ele usa uma descoberta de
beans simplificada e tenta remover beans, interceptors e decorators não
usados durante o build. Isso funciona porque as anotações estão nas
classes que o Quarkus analisa. Uma leitura da Dependency Rule para um
projeto Quarkus: anotações de framework (\texttt{@ApplicationScoped},
\texttt{@Path}, \texttt{@Inject}) ficam nos adaptadores e na camada de
aplicação; as regras de negócio continuam sendo classes e records
comuns, testáveis com JUnit puro, como na Parte~\ref{parte:testes}. Essa divisão é
uma aplicação do livro, não uma regra do Quarkus.

O mesmo guia traz uma recomendação que só faz sentido sabendo como a JVM
acessa membros: no Quarkus, se a injeção de dependência precisa acessar
um membro \texttt{private}, tem de usar reflexão, e em executável nativo
todo membro usado por reflexão precisa ser registrado. Por isso o guia
desencoraja membros \texttt{private} em beans e sugere campos de
injeção \emph{package-private} ou injeção pelo construtor.

\section{Onde aprofundar}

\subsection*{Livros}
\small
\begin{longtable}{@{}p{1.8cm} p{4.4cm} p{2.6cm} p{4cm}@{}}
\toprule
Tema & Livro & Autor(es), ano & Por que ler \\
\midrule \endhead
Conduta & \citetitle{livro:martin2011} \pt{} \emph{O Codificador Limpo}
  & \citeauthor{livro:martin2011}, \citeyear{livro:martin2011}
  & Compromissos, estimativa, pressão e prática \\
Código & \citetitle{livro:martin2008} \pt{} \emph{Código Limpo}
  & \citeauthor{livro:martin2008}, \citeyear{livro:martin2008}
  & Nomes, funções, erros, classes e o catálogo de smells \\
Arquitetura & \citetitle{livro:martin2017} \pt{} \emph{Arquitetura Limpa}
  & \citeauthor{livro:martin2017}, \citeyear{livro:martin2017}
  & Dependency Rule e o framework como detalhe (caps.~22 e 32) \\
Java & \citetitle{livro:bloch2017}, 3ª ed. \pt{} \emph{Java Efetivo}
  & \citeauthor{livro:bloch2017}, \citeyear{livro:bloch2017}
  & Os itens da tabela deste capítulo \\
\bottomrule
\end{longtable}
\normalsize

\subsection*{Documentação}
\citetitle{doc:quarkus-cdi-reference}, seção \enquote{Native Executables
and Private Members}.

\subsection*{De quem fez o Java}
Nomes, bios e mais obras de cada pessoa no apêndice~\ref{ap:pessoas}.
\begin{itemize}
\obra*{video:bloch-api-design}{Design de API com exemplos bons e ruins do
  próprio JDK; vale para qualquer interface pública, inclusive portas da
  arquitetura hexagonal.}
\obra*{artigo:bumper-sticker-api}{As máximas da palestra em uma página.}
\end{itemize}

\subsection*{No IDE: blog da JetBrains}
Outros posts da JetBrains no apêndice~\ref{ap:jetbrains}.
\begin{itemize}
\obra*{artigo:jetbrains-code-smells-mutation}{Mutabilidade como code smell,
  e como records e imutabilidade resolvem.}
\obra*{artigo:jetbrains-code-smells-null}{O \texttt{null} como smell e as
  alternativas.}
\obra*{curso:jetbrains-guide-refactoring}{Refatorar para entender o
  código, com o IDE fazendo o trabalho mecânico.}
\end{itemize}

\subsection*{Pesquisa acadêmica}
Acesso (aberto ou pago) e links no apêndice~\ref{ap:pesquisa}.
\begin{itemize}
\pesquisa{paper:parnas-modules-1972}{Seis páginas que fundaram o
  information hiding: dividir o sistema pelas decisões que podem mudar.}
\pesquisa{paper:silva-why-we-refactor-2016}{Por que desenvolvedores
  refatoram, segundo eles mesmos, logo depois de refatorar.}
\pesquisa{paper:aniche-tese-usp-2016}{Tese do IME-USP: métricas e smells
  dependem do papel da classe (Controller, Repository).}
\pesquisa*{livro:valente2020}{Capítulo~5, sobre princípios de projeto,
  com leitura gratuita no site do livro.}
\pesquisa{doc:stanford-cs190}{O curso de design de software de Stanford que
  deu origem a \emph{A Philosophy of Software Design}.}
\end{itemize}

\section{Desafios práticos}

\desafio{Pegue a última estimativa que você deu no trabalho. Reescreva-a
como três números (otimista, provável, pessimista) e diga o que você
precisaria saber para estreitar o intervalo.}

\desafio{Ache no seu projeto uma afirmação de desempenho que ninguém
mediu (um comentário, uma regra do time, uma \enquote{otimização}).
Descreva que dado confirmaria ou derrubaria a afirmação e qual ferramenta
do capítulo~\ref{cap:lendo-bytecode} você usaria.}

\desafio{Desenhe as dependências entre pacotes de um serviço Quarkus seu.
Alguma regra de negócio importa \texttt{jakarta.*} ou
\texttt{io.quarkus.*}? O que seria preciso para testá-la sem subir o
Quarkus?}

\desafio{Procure campos \texttt{private} com \texttt{@Inject} no seu
projeto Quarkus. Leia a seção \enquote{Native Executables and Private
Members} do guia de CDI e decida, para cada caso, entre campo
package-private e injeção por construtor.}

